\chapter{Происхождение калибровочно-замкнутого рёберно-ходжева запуска}
\label{ch:v8-gauge-closed-edge-hodge-origin-gate}

\section{Канонические операторы полного модуля переноса}

Предыдущий гейт выделил разложение
\begin{equation}
 \mathcal T_{15}=\mathcal I_5\oplus\mathcal H_{10}
 \label{eq:v8-edge-origin-desired-splitting}
\end{equation}
и показал существование устойчивого двухмассового завершения. Теперь нужно
проверить, выводят ли отношение масс уже имеющиеся канонические данные:
полный калибровочный оператор Казимира, кварк-лептонная градуировка, цепная степень и общий
KMS-след.

На представлении переноса полный калибровочный оператор Казимира имеет спектр
\begin{equation}
 \operatorname{Spec}\mathcal C_G
 =\{0^{\times1},1^{\times8},(16/9)^{\times6}\}.
 \label{eq:v8-transfer-gauge-casimir-spectrum}
\end{equation}
Гладкий секторный оператор
\begin{equation}
 \mathcal B(A)=\Gamma_tA-A\Gamma_s
 \label{eq:v8-transfer-sector-order-operator}
\end{equation}
имеет спектр
\begin{equation}
 \operatorname{Spec}\mathcal B
 =\{(-2)^{\times3},0^{\times9},2^{\times3}\}.
 \label{eq:v8-transfer-sector-order-spectrum}
\end{equation}
Операторы $\mathcal C_G$ и $\mathcal B$ коммутируют. Цепной оператор даёт
всем стрелкам переноса одну боровскую степень 2 и потому добавляет только
скалярный оператор.

\section{Совместный спектр и неустранимая кратность}

Совместные спектральные блоки $(\mathcal C_G,\mathcal B)$ имеют комплексные
ранги
\begin{equation}
 (0,0):1,qquad (1,0):8,qquad
 (16/9,-2):3,qquad (16/9,2):3.
 \label{eq:v8-casimir-sector-joint-blocks}
\end{equation}
Проектор инцидентности $P_I$ распределён по ним как
\begin{equation}
 1,qquad 4,qquad0,qquad0.
 \label{eq:v8-incidence-overlap-joint-blocks}
\end{equation}
Иными словами, восьмимерный блок $(1,0)$ содержит четыре инцидентностных и четыре
тяжёлые направления с одинаковыми значениями всех проверяемых операторов.

Наилучшая функция совместного спектра приближает $P_I$ скаляром $1/2$ на
этом блоке. Неустранимые ошибки равны
\begin{equation}
 \|P_I-P_I^{\rm spec}\|_{\rm HS}=\sqrt2,qquad
 \|P_I-P_I^{\rm spec}\|_{\rm op}=\frac12.
 \label{eq:v8-incidence-projector-spectral-residual}
\end{equation}
Для использованной ранее пары $(m_I,m_H)=(4,3.6)$ лучшая спектральная
аппроксимация ставит на вырожденном блоке вес $3.8$ и оставляет ошибку
\begin{equation}
 0.4\sqrt2=0.565685\ldots.
 \label{eq:v8-two-mass-spectral-residual}
\end{equation}

Размерность полного калибровочный коммутанта равна 13 и не уменьшается после
добавления $\mathcal C_G$, $\mathcal B$ и цепной степени. Это показывает,
что они принадлежат уже существующей инвариантной алгебре, но не выбирают
нужную копию внутри кратностного блока.

\section{Почему общий след также не различает копии}

Для центральной концевой плотности
\begin{equation}
 \rho=\rho_sI_{11}\oplus\rho_tI_{10}
 \end{equation}
метрика переноса имеет вид
\begin{equation}
 \langle A,B\rangle_\rho
 = (\rho_s+\rho_t)\Tr(A^*B).
 \label{eq:v8-common-kms-transfer-metric}
\end{equation}
Она скалярна на всём $\mathcal T_{15}$ при след-, KMS- и полярно выведенных
отношениях плотностей концевых секторов. Поэтому общий KMS-след даёт отношение
$m_I/m_H=1$, а не выводит требуемый двухмассовый оператор.

Даже разрешение знака $\mathcal B$ допускает только суммы совместных блоков
комплексных рангов $1,8,3,3$. Ни одна такая сумма не имеет комплексный ранг
5. Следовательно, десятимодовый инцидентностный запуск нельзя получить
спектральным порогом этих операторов.

\begin{theorem}[Кратностный барьер рёберно-ходжева метрики]
Калибровочный оператор Казимира, кварк-лептонная градуировка, общая цепная степень и
центральный KMS-след не восстанавливают проектор $P_I$ и не фиксируют
отношение $m_I/m_H$. Их совместный спектр сохраняет комплексный блок ранга
восемь, содержащий четыре неразличимые инцидентностные и четыре тяжёлые копии.
\end{theorem}

\begin{nogocriterion}[Запрет кругового проектор инцидентностиа]
Нельзя вставлять $P_I$, построенный как калибровочная орбита уже выбранного фона
$A_0$, в исходный гессиан и считать этим происхождение запуска доказанным.
Такое действие использует результат перехода для выбора направлений самого
перехода. Нужен независимый оператор вещественной структуры, бимодульный или мультиплетный оператор,
который расщепляет блок $4+4$ до фиксации вакуума.
\end{nogocriterion}

\begin{protocol}
\begin{enumerate}
 \item спектр калибровочный оператор Казимира: \textbf{$1+8+6$};
 \item совместные блоки $(\mathcal C_G,\mathcal B)$:
 \textbf{$1+8+3+3$};
 \item вырожденный блок: \textbf{$4$ инцидентностные $+4$ тяжёлые копии};
 \item ошибка восстановления $P_I$: \textbf{$\sqrt2$};
 \item ошибка двухмассового представителя: \textbf{$0.4\sqrt2$};
 \item цепная степень различает копии: \textbf{нет};
 \item общий KMS-след различает копии: \textbf{нет};
 \item калибровочный коммутант после всех операторов: \textbf{$13$};
 \item десятимодовый спектральный порог: \textbf{невозможен};
 \item единственная рёберно-ходжева метрика: \textbf{не получена};
 \item полный родительское действие: \textbf{не получен};
 \item следующий гейт: \textbf{вещественно-структурное кратностное факторпространство блока инцидентности}.
\end{enumerate}
\end{protocol}

Машинный сертификат сохранён в
\path{s2t_v8_gauge_closed_edge_hodge_origin_gate_results.json}.